\section{Euler--Boole truncations at every endpoint}\label{sec:boole}
The Euler--Boole method is classical; the point here is the uniform
arithmetic conclusion for this particular family. Define integers
$E_{2j}$ by $\sech t=\sum_{j\geq0}E_{2j}t^{2j}/(2j)!$ near zero.
Multiplication by $\cosh t$ gives their integral recurrence. For $N\geq1$
and $m\geq0$ put
\begin{align}
 S_N&=\sum_{k<N}\frac{(-1)^k}{(2k+1)^4},\notag\\
 T_{N,m}&=\frac1{12}\sum_{j<m}
 \frac{E_{2j}(2j+1)(2j+2)(2j+3)}{(2N)^{2j+4}},\notag\\
 A_{N,m}&=S_N+(-1)^NT_{N,m}=a_{N,m}/q_{N,m},\label{eq:boole-def}
\end{align}
in lowest terms with $q_{N,m}>0$. The Euler-number convention and
classical summation formula are recorded in DLMF~\cite[Sections 24.2,
24.17]{DLMF}; the needed remainder is derived below.

\begin{theorem}[Uniform primitive obstruction; S04]\label{thm:boole}
Every $q_{N,m}\theta-a_{N,m}$ is nonzero. Uniformly in $m\geq1$,
\begin{equation}\label{eq:boole-error}
 |\theta-A_{N,m}|\geq\frac{\pi^2e^{-\pi N}}{144N}.
\end{equation}
Moreover
\[
 \liminf_{N\to\infty}\inf_{m\geq0}\frac{\log q_{N,m}}N\geq8,
 \qquad
 \liminf_{N\to\infty}\inf_{m\geq0}
 \frac{\log|q_{N,m}\theta-a_{N,m}|}N\geq8-\pi>0.
\]
The primitive absolute values tend to infinity whenever $N+m\to\infty$.
\end{theorem}
\begin{proof}
Laplace integration of the alternating tail gives
\begin{equation}\label{eq:boole-tail}
 \theta-S_N=\frac{(-1)^N}{12}
       \int_0^\infty t^3e^{-2Nt}\sech t\dd t.
\end{equation}
Let $a=\pi/2$ and $a_k=(2k+1)a$. The partial fractions formula
\[
 \sech t=2\sum_{k\geq0}\frac{(-1)^k a_k}{a_k^2+t^2}
\]
follows by expanding $\cosh(tx)/\cosh t$ on $[0,1]$ in the mixed
cosine basis $\cos(a_kx)$: twice integrating its differential equation
gives coefficient $2a_k(-1)^k/(t^2+a_k^2)$, and Fourier convergence
at $x=0$ gives the formula.
Applying the finite geometric identity to each term gives, for $m\geq1$,
\begin{align}
 \sech t-\sum_{j<m}\frac{E_{2j}t^{2j}}{(2j)!}
   &=(-1)^m t^{2m}C_m(t),\notag\\
 C_m(t)&=2\sum_{k\geq0}
               \frac{(-1)^k a_k^{1-2m}}{t^2+a_k^2}.\label{eq:boole-rem}
\end{align}
This remainder series is absolutely convergent. Its positive terms
decrease, and the ratio of the second to the first is at most
$3^{1-2m}\leq1/3$. Hence
\begin{equation}\label{eq:boole-sandwich}
 (1-3^{1-2m})\frac{2a^{1-2m}}{t^2+a^2}
 \leq C_m(t)\leq\frac{2a^{1-2m}}{t^2+a^2},\qquad C_m(t)>0.
\end{equation}
Thus $\theta-A_{N,m}=(-1)^{N+m}\epsilon_{N,m}$ with
$\epsilon_{N,m}=\int_0^\infty t^{2m+3}e^{-2Nt}C_m(t)\dd t/12>0$.
On $t\geq a$, use $(t/a)^{2m}\geq1$ and
$t^3/(t^2+a^2)\geq a/2$. The lower bound becomes
$\epsilon_{N,m}\geq(\pi/18)(a/2)e^{-2Na}/(2N)$, proving
\eqref{eq:boole-error}. For $m=0$, $\sech t\geq e^{-t}$ gives
$|\theta-S_N|\geq1/[2(2N+1)^4]$.

For the arithmetic, let $p$ be odd, $p^2>2N$ and $p<2N$, and let
$J$ count the positive odd $r$ with $pr<2N$. Reduction yields
\begin{equation}\label{eq:boole-residue}
 p^4S_N\equiv\chi_4(p)S_J\pmod p,\qquad
 \chi_4(p)=(-1)^{(p-1)/2}.
\end{equation}
If $p\nmid6N$, the correction $T_{N,m}$ is $p$-integral for every
$m$. Whenever $S_J$ is a nonzero residue, therefore,
$v_p(q_{N,m})=4$ exactly. Fix $K\geq1$. For all sufficiently large
$N$, primes in $(2N/(2K+1),2N)$ have $1\leq J\leq K$, satisfy
$p^2>2N$, and avoid the finitely many numerator and denominator prime
divisors of the positive rationals $S_1,\ldots,S_K$. Uniformly in $m$,
\[
 \log q_{N,m}\geq4\{\vartheta(2N)-\vartheta(2N/(2K+1))
                           -\log(6N)\}-O_K(\log N).
\]
First take the lower limit in $N$, then let $K\to\infty$; no uniformity
in $K$ is required. This gives the rate $8$, and combining with the
error estimates gives $8-\pi$.

For bounded $N$ and $m\to\infty$, restrict the positive remainder
integral to $[b,b+1]$, where $b>a$. Its lower bound is
$c_{N,b}(b/a)^{2m}\to\infty$, uniformly over a finite set of $N$.
Since $q\geq1$, this also proves divergence on every escaping
parameter sequence.
\end{proof}

The raw exponential constant is sharp: for $m=\pi N/2+O(1)$,
\begin{equation}\label{eq:boole-sharp}
 |\theta-A_{N,m}|\sim
       \frac{\pi^3}{48\sqrt{2N}}e^{-\pi N}.
\end{equation}
Indeed~\eqref{eq:boole-sandwich} reduces the integral to its first
pole with relative error $O(3^{1-2m})$. The exponent is
$N(2a\log(t/a)-2t)$, with unique maximum at $t=a$, second derivative
$-2/a$, and remaining factor $(\pi/12)t^3/(t^2+a^2)$ there equal
to $\pi a/24$. Gaussian localization gives~\eqref{eq:boole-sharp},
uniformly for the bounded rounding error in $m$.

There are finitely many small primitive values, for example those of
$X-1$, $81X-80$ and $32X-31$. They do not contradict divergence.
Nonvanishing and divergence off finite parameter sets give
$c:=\inf_{N,m}|q_{N,m}\theta-a_{N,m}|>0$. By Gauss's lemma a product
of $d$ of these primitive linear polynomials is primitive, and its
$d$th-root absolute value is at least $c$. Products alone therefore
cannot supply the degree-normalized smallness criterion. Sums,
resultants and general determinant combinations are outside this result.
